% Generated draft: requires mathematical and source review.
\documentclass[11pt]{article}
\usepackage[margin=25mm]{geometry}
\usepackage{fontspec}
\setmainfont{Latin Modern Roman}
\usepackage{amsmath,amssymb,mathtools,unicode-math}
\setmathfont{Latin Modern Math}
\usepackage{enumitem,needspace,xurl,hyperref}
\setlength{\parindent}{0pt}
\setlength{\parskip}{5pt}
\setlength{\emergencystretch}{3em}
\newcommand{\N}{\mathbb{N}}
\newcommand{\Z}{\mathbb{Z}}
\newcommand{\Q}{\mathbb{Q}}
\newcommand{\R}{\mathbb{R}}
\newcommand{\C}{\mathbb{C}}
\newcommand{\F}{\mathbb{F}}
\newcommand{\A}{\mathbb{A}}
\newcommand{\PP}{\mathbb{P}}
\newcommand{\E}{\mathbb{E}}
\newcommand{\eps}{\varepsilon}
\newcommand{\dd}{\,\mathrm{d}}
\begin{document}
\section*{UnsolvedMath ID: 20002005}
\textbf{Source problem code:} catalogue entry
\subsection{Definitions and mathematical statement}
Zharkov's A.23 contains three related research questions, using Gross's tropical Calabi--Yau setup: a compact base \(B\) with an integral affine structure on its nonsingular part \(B_0\), and local potentials differing by affine functions. First, construct natural compact tropicalizations, modifying the valuation map by removing ``pseudo-pods'' while retaining as much affine structure as possible. A polarization may enter the construction, but the desired underlying affine structure should be algebraic and independent of that choice. Second, define a tropical Monge--Amp\`ere operator in a polarization class in \(H^1(B,\operatorname{Aff}(B,\R))\), where the sheaf represents the relevant local affine functions, and determine appropriate existence and uniqueness statements for its measures/potentials. The source proposes piecewise-linear differential graphs in place of smooth \(\dd K\subset\R^n\times(\R^n)^*\), interpreting the primal/dual volume condition distributionally. It warns that lattice-regulated bendings may destroy uniqueness, citing different Voronoi decompositions on higher-dimensional tori.

Third, define tropical Gromov--Witten calculations on Calabi--Yau spaces and determine which tropical curves lift to holomorphic curves. The source describes interval area by pairing its primitive tangent with dual displacement and an eigen-direction restriction for ends at singularities. The singularity model and precise correspondence hypotheses are not specified. Its reference to Mikhalkin's surface correspondence is not an unrestricted theorem for every tropical Calabi--Yau surface; the source itself warns that higher-dimensional tropical curves need not be realizable.

The supplied aid reports only a precise one-dimensional model. Let \(B=\R/L\Z\), \(L>0\), and \(a>0\). Consider convex functions \(K:\R\to\R\) satisfying
\[
K(x+L)-K(x)=ax+b
\]
for some \(b\in\R\), modulo addition of arbitrary affine functions. Distributional differentiation gives a periodic positive Radon measure \(\mu=K''\), descending to \(B\), with exact normalization \(\mu(B)=a\). The aid reports a bijection between these affine-equivalence classes and positive Radon measures of mass \(a\). Thus a prescribed measure determines the potential modulo affine functions; the polarization degree \(a\) alone does not determine a unique measure.

Finite piecewise-linear potentials modulo translation periods correspond exactly to measures with finitely many atoms. To state the integral-slope qualification with \(b\) fixed, choose a cut with no atom at zero and write
\[
\mu=\sum_i m_i\delta_{t_i},\quad0<t_i<L,\quad m_i>0,\quad\sum_i m_i=a.
\]
A potential with the specified \(b\) has integer slopes exactly when
\[
m_i\in\Z\quad\forall i,\qquad
\frac{b-\sum_i m_i(L-t_i)}L\in\Z.
\]
The fraction is its initial slope. If arbitrary real affine functions are quotiented out, \(b\) changes by arbitrary multiples \(cL\); the fixed-\(b\) congruence must therefore not be imposed as an intrinsic condition on the measure alone. This distinguishes real potentials, integral representatives and a fixed polarization normalization.

The aid mentions Berkovich analytification and extended tropicalizations as modern compact replacements, hypothesis-dependent tropical/nonarchimedean Monge--Amp\`ere results, and continuing global obstructions for superabundant curve realizability. It supplies no full references and does not claim a general solution of the three-part program. The appended Chapters B and C begin new material and are not part of A.23.
\subsection{Short English statement}
The source asks about compact tropicalization, tropical Monge--Amp\`ere theory and realizability of tropical curves. The supplied aid gives a circle model where prescribed measures determine convex potentials modulo affine functions. It does not give uniqueness of a measure from a polarization alone or solve the general curve-counting problem.
\subsection{Sources}
\begin{itemize}
\item Ilia Zharkov, A.23, \emph{Compact tropical varieties, Monge--Amp\`ere equation, Calabi conjecture and curve counting}, AIM workshop \emph{Amoebas and tropical geometry}: \url{https://aimath.org/WWN/amoebas/amoebas.pdf}.
\item The source invokes Gross's tropical Calabi--Yau/line-bundle definitions and Mikhalkin's curve-counting work. The aid mentions Berkovich analytification and later tropical/nonarchimedean Monge--Amp\`ere developments without full references. The original literature field is empty.
\item Local registry \texttt{data/unsolvedmath/problems.json}, permanent ID 20002005, code \texttt{catalogue entry}; canonical \texttt{aim-geometry-notes.json}, \texttt{notes[342]}; updated 14 August 2026. Statement status corrected verified records prior PDF comparison and excludes the trailing Chapters B and C. The original ``deleting pseudo-pods'' wording is retained despite corruption in the clean field.
\item Attempt 1: partial result, partially solved status, rigor label proved, novelty confidence low. The machine-generated aid is not peer reviewed and requires independent verification before citation.
\end{itemize}
\end{document}
